% Supplemental draft; does not modify the formal manuscript.
\subsection{Object-level interval correction under geometric mismatch}
\label{supp:object-conformal}

\paragraph{Scope and independent units.}
The conditional depth model can assign narrow intervals despite geometric model error. We evaluated a standard split-conformal correction~\cite{Angelopoulos_2023} as a reporting diagnostic, keeping the unreweighted neutral B8 predictor (structural temperature $T=1$), its $\sigma=0.01$ depth distribution, and its planning policy fixed. The new set contains 48 independently drawn procedural objects: 24 for calibration and 24 for a single held-out evaluation. Both splits follow the same frozen mixture of closed and open structures, regular panels and declared part mismatches, with no rejection or redraw. Each object provides front, front-and-side, around, and occluded observations, giving 192 dependent cases and 576 frames in total. Rendering uses exact poses, noiseless first-hit depth and ideal visible instance/category labels. The target for each case is the true object's extent projected onto the inferred depth axis; axis errors remain part of the evaluation. Generator truth is unavailable to inference. A separate pair with identical acquired fronts and different hidden depths produces exactly equal particle banks, likelihoods and posteriors; this fixture adds no independent object to the test set.

\paragraph{Frozen correction.}
Let $[L_{ic},U_{ic}]$ be the raw 95\% interval for object $i$ under condition $c$, and $y_{ic}$ its projected target extent. The scale floor $s_0=0.04756$\,m is the median, over the original 12 training objects, of their median absolute posterior-median error across four conditions. It uses neither the new calibration objects nor the test objects. We compute
\[
s_{ic}=\max\!\left\{\frac{U_{ic}-L_{ic}}{2},s_0\right\},\qquad
R_i=\max_c\frac{\max\{L_{ic}-y_{ic},\,y_{ic}-U_{ic},\,0\}}{s_{ic}}.
\]
For $n=24$ calibration objects and $\alpha=0.05$, the order statistic at $\lceil(n+1)(1-\alpha)\rceil=24$ gives $q=10.2982$. The corrected set is
\[
\left[\max\{0,L_{ic}-q s_{ic}\},\ U_{ic}+q s_{ic}\right].
\]
The upper endpoint is not clipped to the original 1.4\,m prior support. This standard object-group construction targets simultaneous coverage of the four condition-specific extents of a new object, marginally over exchangeable objects from the same mixture. It does not assert coverage conditional on a particular class or observation direction. The interval expansion changes neither the probability mass function nor candidate scores. CRPS therefore remains a property of the original distribution; interval width and interval score assess the corrected set separately.

\paragraph{Coverage--width tradeoff.}
All 96 held-out cases remain in the denominator, including eight cases with multiple observed tracks that use the predeclared public uniform-prior/world-$y$ fallback. Calibration retains five fallback cases. The original neutral posterior attains CRPS $0.1099$\,m, compared with $0.1887$\,m for its prior-only control and $0.3146$\,m for a fixed 0.65\,m estimate. The correction increases case coverage from $61/96$ to $93/96$ and all-four-condition coverage from $5/24$ to $21/24$ objects (Table~\ref{tab:b9-intervals}). However, mean width grows from $0.278$ to $2.316$\,m, and mean interval score worsens from $2.140$ to $2.704$\,m. All three remaining misses occur in front-only observations. The observed object coverage of $87.5\%$ is below the nominal $95\%$ target; the marginal conformal statement does not guarantee that every finite test batch achieves its nominal fraction. These wide intervals are retained as a diagnostic and are not adopted by the planner. The experiment supplies no evidence that set expansion improves point geometry, probability scores, navigation or reconstruction.

\begin{table}[tbp]
\centering\small
\caption{Unchanged neutral depth distributions on the same 24 test objects (96 dependent cases). All eight evaluator fallbacks remain in the denominator. CRPS and median error are in metres. The set correction in Table~\ref{tab:b9-intervals} does not change these probability scores.}
\label{tab:b9-original-pmf}
\begin{tabular}{lrr}
\toprule
Prediction & CRPS $\downarrow$ & Median error $\downarrow$ \\
\midrule
Frozen neutral posterior & 0.1099 & 0.1388 \\
Conditional prior only & 0.1887 & 0.2882 \\
Fixed 0.65 m & 0.3146 & 0.3146 \\
\bottomrule
\end{tabular}
\end{table}

\begin{table}[tbp]
\centering\small
\caption{Report-only 95\% interval correction on 24 independent test objects, each with four dependent observations. Calibration uses 24 separate objects and the frozen neutral depth predictor. Case coverage and all-four-condition object coverage have distinct denominators. Width and interval score are in metres; lower values are better. The corrected intervals leave the original probability distribution and navigation unchanged.}
\label{tab:b9-intervals}
\begin{tabular}{llrrrr}
\toprule
Condition & Interval & Cases covered & Objects covered & Width $\downarrow$ & Score $\downarrow$ \\
\midrule
All conditions & Raw & 61/96 & 5/24 & 0.278 & 2.140 \\
All conditions & Corrected & 93/96 & 21/24 & 2.316 & 2.704 \\
\midrule
Front & Raw & 13/24 & -- & 0.440 & 6.652 \\
Front & Corrected & 21/24 & -- & 3.242 & 4.794 \\
Front + side & Raw & 19/24 & -- & 0.131 & 0.444 \\
Front + side & Corrected & 24/24 & -- & 1.399 & 1.399 \\
Around & Raw & 9/24 & -- & 0.062 & 0.325 \\
Around & Corrected & 24/24 & -- & 1.227 & 1.227 \\
Occluded & Raw & 20/24 & -- & 0.478 & 1.137 \\
Occluded & Corrected & 24/24 & -- & 3.395 & 3.395 \\
\bottomrule
\end{tabular}
\end{table}

